\documentclass[11pt]{article}

\usepackage[english]{babel}
\usepackage[latin1]{inputenc}
\usepackage{graphicx}
\usepackage{amssymb}
\usepackage{url}

%\usepackage{mathtools}
%\mathtoolsset{showonlyrefs}

\usepackage{color}
\usepackage{natbib}
% \usepackage[style=authoryear-comp,dashed=false]{biblatex}
% \bibliography{piecewise_lyapunov}

\textwidth 15cm
\setlength{\textheight}{1.1\textheight}
\newcommand\hi{\hspace*{\parindent}}
\newcommand\vi{\vspace{\baselineskip}}
\newcommand\lac{{\mbox{{\Huge\bf L}\hspace{-0.65em}
\raisebox{-1.2ex}{\Huge\bf A}\hspace{-1.1em}
\raisebox{-0.6ex}{\Huge\bf C~}}}}


\newcommand{\Real}{\mathbb{R}}



\pagestyle{empty}

\begin{document}

\hoffset -1.5cm \voffset .5cm


\sf

\vspace*{-2.5cm} \hspace*{-0.1cm} {\small
{\mbox{\begin{minipage}{1.5cm}
\centerline{\includegraphics[width=2cm]{manchester.jpg}}
\end{minipage}}} \hspace*{0.1cm}
\begin{minipage}{11.0cm}{\sf \large \textbf{UNIVERSITY OF MANCHESTER}}\\
{\sc School of Computer Science}
\end{minipage}
} \vspace*{0mm}

\hspace*{-.7cm} {\rule[-1ex]{15cm}{0.03cm}}

\begin{flushright}
\begin{minipage}{7.0cm}\small
{\bf Please Reply to:}\\
Dr. Cordeiro\\
University of Manchester\\
School of Computer Science\\
Kilburn Building, Manchester M13 9PL\\
{\em e-mail: lucas.cordeiro@manchester.ac.uk}
\end{minipage}
\end{flushright}

\vi
\hspace*{\fill}{\small Manchester, \today.}
\vi

\begin{flushleft}
      Badrul H. Chowdhury,  Ph.D. \\
	Editor-in-Chief \\
	University of North Carolina at Charlotte \\
	IEEE Transactions on Sustainable Energy 
    \end{flushleft}
 
\indent Dear Badrul
\vi

Please find enclosed our manuscript entitled as ``Automated Verification of Stand-alone Solar Photovoltaic Systems'' by Alessandro Trindade and Lucas Cordeiro which we would like to submit for publication in the journal of IEEE Transactions on Sustainable Energy.

The present work describes and evaluates a formal verification methodology for stand-alone solar photovoltaic systems, which employs the Efficient SMT-based Bounded Model Checker (ESBMC) tool to check the most common type of renewable energy used in rural areas or where grid extension unfeasible.

The evaluation of PV systems is usually performed by simulation tools, which are well known in the market. However, the exploration of all design space state is impossible with those tools, and some flaws can be originated from the design phase, reaching the field, after the PV deployment. This can cause dissatisfaction to the PV system owners, and to the wrong conclusion that intermittent renewable systems are not good.
   
In our study, experimental results, from five case studies, in a comparative evaluation of tools, showed that only the automated verification can find some failures in the PV systems, and with the plus of indicating the reason of failure. That can and must be used to improve the sizing of solar PV systems.

Finally, based on the fact that only since 2017 papers are dealing with formal verification applied to PV systems, with excellent results, but only related to PV panels or to grid-tied systems, our study is unique and based on data from four months of use from the case studies: putting together mathematical models, automated verification, simulation, empirical observation and interview from dwellers in remote communities.

The  authors  claim  that  none  of  the  material  in  the mentioned  paper  has  been  published  or  is  under  consideration for publication elsewhere.

\vi

Thank you for your time.
\vi

\indent
~~~~~~~~~~~~~~~~~~~~~~~~~~~Sincerely,\\

\begin{quote}
\begin{quote}
\begin{flushright}


%\vfill

Lucas C. Cordeiro~~~~~~~
\end{flushright}
\end{quote}
\end{quote}


 
 

\end{document} 
